\documentclass[lettersize,journal]{IEEEtran}
\usepackage{amsmath,amsfonts,amssymb}
\usepackage{algorithmic}
\usepackage{algorithm}
\usepackage{array}
%\usepackage[caption=false,font=normalsize,labelfont=sf,textfont=sf]{subfig}
\usepackage[caption=false,font=normalsize]{subfig}
\usepackage{textcomp}
\usepackage{stfloats}
\usepackage{url}
\usepackage{verbatim}
\usepackage{graphicx}
\usepackage{cite}
\usepackage{hyperref}
\usepackage{bm}
\usepackage{amsthm}
\newtheorem{lemma}{Lemma}
\newtheorem{proposition}{Proposition}
\newtheorem{corollary}{Corollary}
\theoremstyle{remark}
\newtheorem{remark}{Remark}
% \usepackage[hidelinks]{hyperref}
\usepackage{xcolor}
\newcommand{\rev}[1]{\textcolor{blue}{#1}}      % mark revisions
\newcommand{\todo}[1]{\textcolor{red}{[TODO: #1]}} % inline TODO marker
\hyphenation{op-tical net-works semi-conduc-tor IEEE-Xplore}

\begin{document}

\title{Projection-Retraction MPPI: \\ Exact Constraint-Manifold Control for Manipulators}

% Author intentionally omitted for blind first submission.
% For the camera-ready, add:
% \author{Author Name(s)\thanks{Affiliation / funding note.}}

\maketitle

\begin{abstract}
In many manipulation tasks a robot must hold a hard kinematic equality constraint throughout the
\end{abstract}

\begin{IEEEkeywords}
Sampling-based model predictive control, model predictive path integral, equality-constrained
control, null-space projection, dual-arm manipulation, constrained motion control.
\end{IEEEkeywords}

%==============================================================
\section{INTRODUCTION}
\label{sec:intro}
%==============================================================
\IEEEPARstart{M}{any} manipulation tasks ask the robot to do more than reach a goal: they require
it to hold a kinematic equality constraint on its configuration throughout the motion. A wiping or
polishing tool must stay on the work surface; a carried tray or glass must stay level; two arms
that rigidly hold the same object form a closed kinematic chain whose relative pose must not
change. Each case ties the joint configuration to a fixed relationship that must hold
\emph{exactly} rather than on average---a closed chain that drifts even slightly
builds up internal wrenches or loses the grasp, and a tray that tilts a few degrees spills. The
controller must therefore keep the robot on the constraint at every step while it makes progress on
the task.


%==============================================================
\section{CONCLUSION}
\label{sec:conclusion}
%==============================================================
test\cite{bishop2006pattern}

\bibliographystyle{IEEEtran}
\bibliography{IEEEabrv, references, ours}

\end{document}
